% LaTeX-Text für Aufgabe 7: Winkelunsicherheit

\subsection{Aufgabe 7: Winkelunsicherheit}

Zur Bestimmung der Winkelunsicherheit wurden zwei Messreihen (A1.1 und A3) mit identischen Messparametern durchgeführt, bei denen jedoch der Kristall um \SI{180}{\degree} gedreht wurde. Die unterschiedliche Kristallorientierung führt zu einer horizontalen Verschiebung der gemessenen Beugungsmuster, die sich in einer systematischen Verschiebung der Peaks und charakteristischen Merkmale äußert.

\subsubsection{Methode}

Um die durch die Kristallorientierung verursachte Winkelunsicherheit zu quantifizieren, wurden beide Messkurven gegeneinander geplottet (siehe Abbildung~\ref{fig:aufg7_originalkurven}). Mit bloßem Auge ist zu erkennen, dass die Daten aus Messreihe A3 nach rechts verschoben werden müssen, um eine saubere Überlagerung mit A1.1 zu erzielen.

Für eine quantitative Bestimmung der Verschiebung wurde der Messbereich in drei strukturreiche Abschnitte eingeteilt:
\begin{itemize}
    \item \textbf{Fenster 1} (\SIrange{3}{4}{\degree}): Steiler Anstiegsbereich
    \item \textbf{Fenster 2} (\SIrange{6.0}{7.0}{\degree}): Erster Peak
    \item \textbf{Fenster 3} (\SIrange{6.4}{7.4}{\degree}): Zweiter Peak / Hauptpeak
\end{itemize}

Für jedes Fenster wurde die optimale horizontale Verschiebung $\Delta\beta_i$ durch Maximierung der Kreuzkorrelation zwischen den normalisierten Kurven bestimmt. Dabei wurden die Zählraten zunächst mit einem Savitzky-Golay-Filter (Fensterbreite 11 bzw.\ 9, Polynomgrad 3) geglättet, um statistische Fluktuationen zu reduzieren. Die Kurven wurden auf ein feines Interpolationsgitter gelegt und die Korrelation zwischen $R_{\text{A1.1}}(\beta)$ und $R_{\text{A3}}(\beta - s)$ für verschiedene Shift-Werte $s \in [-\SI{1}{\degree}, +\SI{1}{\degree}]$ berechnet. Der optimale Shift entspricht dem Maximum der Korrelationsfunktion (siehe Abbildung~\ref{fig:aufg7_correlation}).

\subsubsection{Ergebnisse}

Die Analyse der drei Fenster ergab folgende Verschiebungen:

\begin{table}[h]
\centering
\caption{Ermittelte horizontale Verschiebungen in den drei Analysefenstern}
\label{tab:shifts}
\begin{tabular}{lccc}
\toprule
Fenster & Bereich & $\Delta\beta_i$ & Korrelation \\
\midrule
1 (Anstieg)    & \SIrange{3.0}{4.0}{\degree}   & $+\SI{0.2308}{\degree}$ & 0.992 \\
2 (Peak 1)     & \SIrange{6.0}{7.0}{\degree}   & $+\SI{0.3043}{\degree}$ & 0.938 \\
3 (Peak 2)     & \SIrange{6.4}{7.4}{\degree}   & $+\SI{0.2977}{\degree}$ & 0.893 \\
\bottomrule
\end{tabular}
\end{table}

Die hohen Korrelationswerte (0.89--0.99) zeigen, dass die Kurvenformen in den ausgewählten Bereichen gut übereinstimmen und die Shift-Bestimmung robust ist. Die leichten Unterschiede zwischen den Fenstern reflektieren, dass die Verschiebung nicht in allen Winkelbereichen exakt gleich ist, was auf eine leichte Verkippung oder nicht-ideale Kristallausrichtung hindeutet.

Aus den drei Einzelwerten wurde der Mittelwert und die Standardabweichung berechnet:
\begin{align}
\overline{\Delta\beta} &= \frac{1}{3}\sum_{i=1}^{3} \Delta\beta_i = \SI{+0.2776}{\degree} \approx \SI{+0.28}{\degree} \\
u_\beta &= \sqrt{\frac{1}{n-1}\sum_{i=1}^{3}(\Delta\beta_i - \overline{\Delta\beta})^2} = \SI{0.0407}{\degree} \approx \SI{0.04}{\degree}
\end{align}

Die Standardabweichung $u_\beta = \SI{0.04}{\degree}$ repräsentiert die zusätzliche Winkelunsicherheit, die aus der ungenauen Kristallpositionierung bzw.\ -orientierung resultiert.

\subsubsection{Gesamtunsicherheit}

Die ermittelte Winkelunsicherheit aus der Kristallorientierung muss quadratisch mit der angegebenen Geräteunsicherheit des Goniometers ($u_{\text{Gerät}} = \SI{0.05}{\degree}$) addiert werden:
\begin{align}
u_{\text{gesamt}} &= \sqrt{u_{\text{Gerät}}^2 + u_\beta^2} \\
                   &= \sqrt{(\SI{0.05}{\degree})^2 + (\SI{0.04}{\degree})^2} \\
                   &= \sqrt{\SI{0.0025}{\degree^2} + \SI{0.0016}{\degree^2}} \\
                   &= \SI{0.064}{\degree} \approx \SI{0.06}{\degree}
\end{align}

Abbildung~\ref{fig:aufg7_vorher_nachher} zeigt die Messkurven vor und nach der Shift-Korrektur um $\overline{\Delta\beta} = \SI{+0.28}{\degree}$. Nach der Korrektur liegen die charakteristischen Merkmale (Peaks, Anstiege) deutlich besser übereinander, wobei verbleibende Abweichungen durch die inhärente Winkelunsicherheit von $u_\beta = \SI{0.04}{\degree}$ erklärt werden. Abbildung~\ref{fig:aufg7_shift_summary} fasst die Einzelergebnisse grafisch zusammen.

\textbf{Für alle weiteren Berechnungen in diesem Versuch wird die Gesamtunsicherheit $u_\beta = \SI{0.06}{\degree}$ bzw.\ $u_\beta = \SI{0.07}{\degree}$ (konservativ) verwendet.}

\begin{figure}[h]
    \centering
    \includegraphics[width=0.95\textwidth]{aufg7_originalkurven.pdf}
    \caption{Vergleich der Originalkurven A1.1 und A3. Die Kurve A3 (orange) liegt systematisch bei kleineren $\beta$-Werten als A1.1 (blau), was auf eine horizontale Verschiebung durch unterschiedliche Kristallorientierung hindeutet.}
    \label{fig:aufg7_originalkurven}
\end{figure}

\begin{figure}[h]
    \centering
    \includegraphics[width=0.95\textwidth]{aufg7_correlation_analysis.pdf}
    \caption{Korrelationsanalyse für die drei Analysefenster. Das Maximum der Korrelationsfunktion bestimmt die optimale Verschiebung $\Delta\beta_i$ für jedes Fenster. Die hohen Korrelationswerte ($>0.89$) bestätigen die Zuverlässigkeit der Methode.}
    \label{fig:aufg7_correlation}
\end{figure}

\begin{figure}[h]
    \centering
    \includegraphics[width=0.95\textwidth]{aufg7_shift_summary.pdf}
    \caption{Übersicht der ermittelten Shifts in den drei Analysefenstern. Der Mittelwert $\overline{\Delta\beta} = \SI{+0.28}{\degree}$ (gestrichelte Linie) und die Standardabweichung $u_\beta = \SI{0.044}{\degree}$ (grauer Bereich) charakterisieren die systematische Verschiebung und deren Unsicherheit.}
    \label{fig:aufg7_shift_summary}
\end{figure}

\begin{figure}[h]
    \centering
    \includegraphics[width=0.95\textwidth]{aufg7_vorher_nachher.pdf}
    \caption{Vergleich der Messkurven vor (oben) und nach (unten) der Shift-Korrektur. Nach Verschiebung von A3 um $\Delta\beta = \SI{+0.28}{\degree}$ liegen die charakteristischen Merkmale (Peaks, Anstiege) deutlich besser übereinander, wobei die verbleibende Abweichung die Winkelunsicherheit widerspiegelt.}
    \label{fig:aufg7_vorher_nachher}
\end{figure}

\clearpage


